\subsection{高为 1 的赋值}

命题 6. 设 K 是域，A 是 K 的一个子环。假设 A 不是域。那么下列条件等价：

\begin{enumerate}
\def\labelenumi{(\alph{enumi})}
\item
  A 是 K 上某个高为 1 的赋值的环；
\item
  A 是 K 的赋值环，且除 (0) 与 m(A) 外没有其他素理想；
\item
  A 在异于 K 的 K 的子环中是极大的。
\end{enumerate}

第 4 节命题 5 表明 (a) 蕴含 (b)，第 1 节命题 1 表明 (b) 蕴含 (c)。剩下的要证明 (c) 蕴含 (a)。设 A 在异于 K 的 K 的子环中是极大的。设 m 是 A 的一个极大理想，V 是支配 \(A_{m}\) 的 K 的一个赋值环（§ 1 第 2 节定理 2 的推论）；由于 \(m(V) \cap A = m\) 且 \(m \neq (0)\) （因为 A 不是域），\(V \neq K\) ，于是 V = A，这证明 A 是 K 上某个赋值 v 的环。这样，由（第 1 节）命题 1 与（第 4 节）命题 5，v 的高为 1。

命题 7. 域上的一个赋值的高为 1，当且仅当其阶群同构于 R 的一个非零有序子群。

这实际上由下面的命题得出：

命题 8. 设 G 是一个不退缩为 0 的全序群。下列条件等价：

\begin{enumerate}
\def\labelenumi{(\alph{enumi})}
\item
  G 的高为 1；
\item
  对所有 \(x > 0\) 和 \(y \geqslant 0\) （\(G\) 中），存在整数 \(n \geqslant 0\) 使得 \(y \leqslant nx\) ；
\item
  G 同构于有序加法群 R 的一个不退缩为 0 的子群。
\end{enumerate}

设 \(x\) 是 \(G\) 的一个正元素，\(H_x\) 是由这样的 \(y \in G\) 组成的集合：存在整数 \(n \geqslant 0\) 满足 \(|y| \leqslant nx\) 。容易验证 \(H_x\) 是 \(G\) 的孤立子群，并且 \(G\) 的每个含 \(x\) 的孤立子群都包含 \(H_x\) 。于是条件 (a) 等价于“ \(H_x = G\) 对一切 \(x > 0\) 成立”，即等价于条件 (b)。

显然 (c) 蕴含 (b)。反之，设条件 (b) 成立，并用 Q 表示 G 中 >0 的元素组成的集合。先设 Q 有最小元素 x；对一切 \(y \in Q\) ，设 n 是使 \(y \leqslant nx\) 的最小整数；若 y < nx，则还有 \(nx - y \geqslant x\) ，于是 \(y \leqslant (n - 1)x\) ，与 n 的选取矛盾；从而 y = nx，这表明 \(G = Zx\) 同构于 \(Z \subset R\) 。再设 Q 没有最小元素；我们对有序集 \(P = Q \cup \{0\}\) 应用《一般拓扑学》第五章 §2 命题 1（这是可行的，因为条件 (b) 正是“Archimedes 公理”）；可见存在由 P 到 \(R_{+}\) 的严格递增映射 f，使得

\[
f (x + y) = f (x) + f (y)
\]

对 \(x \in P\) 与 \(y \in P\) 成立；由线性性，f 可以延拓为 G 到 R 的某个子群上的同构，这证明 (b) 蕴含 (c)。

命题 9. 设 K 是域，v 是 K 上一个非真的赋值，A 是 v 的环。A 为完全整闭的（第五章 § 1 第 4 节定义 5），当且仅当 v 的高为 1。

设 \(v\) 的高为 1。设 \(x \in \mathbf{K}\) 使得诸 \(x^n (n \geqslant 0)\) 都含于 \(\mathbf{K}\) 的某个有限生成子 A-模中。于是存在 \(d \in \mathbf{A} - \{0\}\) 使得 \(dx^n \in \mathbf{A}\) 对一切 \(n \geqslant 0\) 成立。那么 \(v(d) + nv(x) \geqslant 0\) ，即 \(n(-v(x)) \leqslant v(d)\) 对一切 \(n \geqslant 0\) 成立，于是 \(-v(x) \leqslant 0\) （命题 8 (b)），因而 \(x \in \mathbf{A}\) 。这样 \(\mathbf{A}\) 是完全整闭的。

现在设 \(v\) 的高不是 1。那么存在 \(y \in \mathfrak{m}(\mathbf{A})\) 与 \(t \in \mathbf{A}\) 使得 \(nv(y) < v(t)\) 对一切 \(n \geqslant 0\) 成立（命题 8 (b)）。于是 \(ty^{-n} \in \mathbf{A}\) 对一切 \(n \geqslant 0\) 成立，但 \(y^{-1} \notin \mathbf{A}\) 。因此 \(\mathbf{A}\) 不是完全整闭的。

推论. 设 \(\mathbf{K}\) 是域，\((v_{\alpha})_{\alpha \in \mathbf{I}}\) 是 \(\mathbf{K}\) 上一族高为 1 的赋值，\(\mathbf{A}\) 是诸 \(v_{\alpha}\) 的环的交。那么 \(\mathbf{A}\) 是完全整闭整环。

完全整闭整环不一定总是高为 1 的赋值环的交（习题 6）。

